\documentclass{article}
\usepackage{amsmath, amssymb, amsthm}

\setlength{\parindent}{0pt}

\newtheorem{claim}{Claim}

\begin{document}

Suppose a generalised World Series between the Sox and the Cardinals involved $2n + 1$ games. As 
usual, the generalised Series will stop as soon as one team has won more than half the possible 
games.

\begin{claim}
    If the Cards won exactly $r$ games and the Sox won the GSeries, then there are 
    $\binom{n + r}{r}$ possible win-loss patterns for the Sox.
\end{claim}

\begin{proof}
    By definition of the GSeries, the Sox must win $n + 1$ games to win the series. Since the Cards 
    win exactly $r$ games, the Sox have exactly $r$ losses. Hence, the series lasts for 
    \[
    n + r + 1
    \]
    games.
    
    \medskip
    
    The Sox must win the final game, since this is the game in which they obtain their $(n+1)$-st 
    win and clinch the GSeries. Therefore, the first $n + r$ games contain exactly $r$ Sox losses 
    and $n$ Sox wins.

    \medskip

    Thus, a win-loss pattern is uniquely determined by choosing the $r$ games among the first 
    $n + r$ games in which the Sox lose. The number of possible choices is
    \[
    \binom{n + r}{r}
    \]

    \medskip

    Hence the Sox have $\binom{n + r}{r}$ possible win-loss patterns.
\end{proof}

\begin{claim}
    If the Cards won at most $r$ games and the Sox won the GSeries, then there are 
    $\binom{n + r + 1}{r}$ possible win-loss patterns for the Sox.
\end{claim}

\begin{proof}
    By definition of the GSeries, the Sox must win $n + 1$ games to win the series. Since the Cards 
    win exactly $r$ games, the Sox have exactly $r$ losses. Hence, the series lasts for 
    \[
    n + r + 1
    \]
    games.

    \medskip

    Represent each win-loss pattern by a binary sequence consisting of $n + r + 1$ bits, where a $1$ 
    represents a Sox win and a $0$ represents a Cards win. If the Cards win fewer than $r$ games, 
    append the remaining $0$s after the final $1$. Thus the sequence consists of exactly 
    $n + 1$ $1$s and $r$ $0$s.

    \medskip

    Conversely, given any such binary sequence, the number of $0$s occurring before the final $1$
    is exactly the number of games won by the Cards before the Sox clinch the GSeries. Any $0$s 
    occurring after the final $1$ are padding bits and do not correspond to any games. Deleting 
    all padding $0$s after the final $1$ recovers the original win-loss pattern. Hence this 
    correspondence is one-to-one.

    \medskip

    Thus, the number of possible win-loss patterns is equal to the number of ways to choose the 
    positions of the $r$ $0$s among the $n + r + 1$ bits, which is
    \[
    \binom{n + r + 1}{r}
    \]

    \medskip

    Hence the Sox have $\binom{n + r + 1}{r}$ possible win-loss patterns.
\end{proof}

\begin{claim}
    \[
    \sum_{i=0}^{r}\binom{n + i}{i} = \binom{n + r + 1}{r}
    \]
\end{claim}

\begin{proof}[Combinatorial Proof]
    From Claim 2, we know that $\binom{n + r + 1}{r}$ is the number of possible win-loss patterns 
    for the Sox given that the Cards won at most $r$ games.

    \medskip

    Now, let $i$ denote the number of games won by the Cards, where $0 \leq i \leq r$. By Claim 1, 
    there are $\binom{n + i}{i}$ possible win-loss patterns for the Sox given that the Cards won 
    $i$ games. Summing over all admissible values gives
    \[
    \sum_{i=0}^{r}\binom{n + i}{i}
    \]

    \medskip

    Since both expressions count the number of possible win-loss patterns for the Sox given that 
    the Cards won at most $r$ games, they must be equal.

    \medskip

    Hence
    \[
    \sum_{i=0}^{r}\binom{n + i}{i} = \binom{n + r + 1}{r}
    \]
\end{proof}

\begin{proof}[Inductive Proof]
    Fix $n \ge 0$. We proceed by induction on $r$.

    \medskip

    \textbf{Base case:} $r = 0$, so
    \[
    \begin{aligned}
    \sum_{i=0}^{0}\binom{n + i}{i} &= \binom{n + 0}{0} \\
    &= 1 \\
    &= \binom{n + 1}{0}
    \end{aligned}
    \]
    Thus the base case holds.

    \medskip

    \textbf{Inductive step:} Assume that for some $r \ge 0$, the identity
    \[
    \sum_{i=0}^{r}\binom{n + i}{i} = \binom{n + r + 1}{r}
    \]
    holds. We must show that
    \[
    \sum_{i=0}^{r+1}\binom{n + i}{i} = \binom{n + r + 2}{r + 1}
    \]

    \medskip

    We have
    \[
    \begin{aligned}
    \sum_{i=0}^{r+1}\binom{n + i}{i} &= \sum_{i=0}^{r}\binom{n + i}{i} + \binom{n + r + 1}{r + 1} \\
    &= \binom{n + r + 1}{r} + \binom{n + r + 1}{r + 1} \qquad \text{(by the induction hypothesis)} \\
    &= \binom{n + r + 2}{r + 1} \qquad \text{(by Pascal's identity).}
    \end{aligned}
    \]
    Thus the inductive step holds.

    \medskip

    Hence, by induction, the identity holds for all $r \ge 0$. Since $n$ was arbitrary, it follows that
    \[
    \sum_{i=0}^{r}\binom{n + i}{i} = \binom{n + r + 1}{r}
    \]
    for all $n, r \ge 0$
\end{proof}

\end{document}
